\chapter{The Layered Stack: A Reference Map}
\label{app:stack}

\chapterorient{This book taught one description of the simulator---the $L_4$
calculus of Part~III---grounded in the cited source and pure-function operators of
Parts~II and~IV, and composed into the modalities of Part~V. But the calculus is
only the top of a five-layer stack, and the full specification the book draws on is
built from all five layers and the rotations between them. This appendix is the
reference map of that stack. It is not a tutorial---\cref{ch:rotations} already
taught the rotation idea, and you should read that chapter first. It is a
\emph{legend}: a one-page-per-layer account of what each layer expresses and what
impedance it has turned; a reference table of the four rotations; a catalogue of the
obstructions that do not rotate; and a chart of how the book's parts map onto the
layers. Keep it beside you when you go to read the underlying spec.}

This appendix maps the layered representation that underlies the entire book. Where
\cref{ch:rotations} \emph{taught} the rotation idea---what a rotation is, the three
tests that detect one, the five justification kinds, and how an obstruction is an
honest failure---this appendix \emph{tabulates} the structure for reference. The
reader who has finished the book and now wants to open the full specification will
find here the map: the five layers in order, the four edges between them, the
negative results, and the correspondence to the parts already read.

\begin{keyidea}
The stack is \textbf{five layers, four rotations, and a catalogue of obstructions}.
The five layers ($L_0$ source, $L_1$ pure functions, $L_2$ algebra, $L_3$
tensor-field operations, $L_4$ calculus) are each a complete, correct account of the
simulator in its own vocabulary. The four rotations between adjacent layers each turn
exactly one impedance---mutation, fusion, iteration, coordination---without changing
the numbers. And the obstructions---the loops that will \emph{not} rotate to a
whole-field form---are findings, not failures: they mark precisely where the
abstract tensor picture meets a wall, and why. The book taught $L_4$; this map shows
the rest.
\end{keyidea}

\section{The five layers}
\label{sec:fivelayers}

We number the layers from the ground up: $L_0$ is the literal source, $L_4$ the
calculus. \Cref{fig:appstack} is the hero figure of this appendix---the polished
reference version of the stack first drawn in \cref{ch:rotations}. It shows all five
bands and the four labeled rotation arrows between them, each tagged with the
impedance it turns. The subsections that follow give each layer its own paragraph:
what it expresses, and what impedance it has rotated relative to the layer below.

\begin{figure}[p]
\centering
\begin{tikzpicture}[node distance=0mm]
  \def\lw{96mm}    % layer width
  \def\lh{14mm}    % layer height
  \def\gap{12mm}   % vertical gap between bands (room for rotation arrow + 2 labels)
  % five stacked layer bands, L0 at bottom -> L4 at top
  \foreach \i/\name/\desc/\col/\bg in {%
      0/$L_0$/the cited source {\textbullet} Palace C++ {\&} MFEM {\textbullet} reference notes/simInk/simBgGray,
      1/$L_1$/pure functions over immutable tensors {\textbullet} no in-place mutation/simGreen/simBgGreen,
      2/$L_2$/algebraic decompositions {\textbullet} HPC fusions erased/simTeal/simBgGray,
      3/$L_3$/global tensor-field operations {\textbullet} element loops lifted/simBlue/simBgBlue,
      4/$L_4$/the calculus {\textbullet} combinators {\textbullet} state monad {\textbullet} strata/simViolet/simBgViolet}{
    \node[draw=\col, fill=\bg, very thick, rounded corners=3pt,
          minimum width=\lw, minimum height=\lh, align=center, font=\small]
      (L\i) at (0, \i*\dimexpr\lh+\gap\relax) {%
        \textbf{\name}\quad\textcolor{simGray}{\footnotesize\desc}};
  }
  % rotation arrows between adjacent bands
  \foreach \lo/\hi in {0/1,1/2,2/3,3/4}{
    \draw[depedge, line width=1.3pt] (L\lo.north) -- (L\hi.south);
  }
  % rotation labels (name + impedance turned), placed to the RIGHT of each arrow gap
  \def\lx{0.52*\lw+3mm}
  \foreach \i/\rot/\imp in {%
      0/mutation rotation/{mutation $\to$ purity},
      1/fusion rotation/{fused kernel $\to$ base algebra},
      2/iteration rotation/{element loop $\to$ whole-field op},
      3/calculus rotation/{control flow $\to$ named combinators}}{
    \node[font=\scriptsize\bfseries\itshape, simRust, anchor=west]
      at (\lx, \i*\dimexpr\lh+\gap\relax+0.5*\lh+0.5*\gap+1.6mm) {\rot};
    \node[font=\scriptsize, simGray, anchor=west]
      at (\lx, \i*\dimexpr\lh+\gap\relax+0.5*\lh+0.5*\gap-1.6mm) {turns: \imp};
  }
  % side bracket: "complete & correct in itself"
  \draw[decorate, decoration={brace, amplitude=6pt}, simGray, thick]
    (-0.52*\lw-3mm, -0.5*\lh) -- (-0.52*\lw-3mm, 4*\dimexpr\lh+\gap\relax+0.5*\lh);
  \node[rotate=90, font=\scriptsize\itshape, simGray, anchor=south]
    at (-0.52*\lw-9mm, 2*\dimexpr\lh+\gap\relax) {each layer complete \& correct in its own vocabulary};
  % top / bottom annotations
  \node[font=\scriptsize\itshape, simViolet, anchor=south] at (0, 4*\dimexpr\lh+\gap\relax+0.5*\lh+2mm)
    {\textbf{top:} the accelerator-ready surface (what the backend consumes)};
  \node[font=\scriptsize\itshape, simInk, anchor=north] at (0, -0.5*\lh-2mm)
    {\textbf{ground:} every claim in the stack rests on an $L_0$ citation};
\end{tikzpicture}
\caption[The five-layer impedance-matching stack]{The impedance-matching stack, in
full. Read it from the ground up. $L_0$ is the cited Palace/MFEM source---the only
layer that \emph{is} the machine rather than a description of it. Each rotation above
turns exactly one impedance while computing the same numbers: \emph{mutation} into
purity ($L_0\!\to\!L_1$), a \emph{fused kernel} into the base algebra it implements
($L_1\!\to\!L_2$), an explicit \emph{element loop} into a whole-field operation
($L_2\!\to\!L_3$), and hand-written \emph{control flow} into named combinators
($L_3\!\to\!L_4$). Every band is a complete, correct account of the simulator in its
own vocabulary; the arrows are translations between vocabularies, never mere
renamings (\cref{ch:rotations}). The book has been teaching the $L_4$ view at the top,
with Parts~II and~IV supplying the $L_0$--$L_1$ grounding at the bottom.}
\label{fig:appstack}
\end{figure}

\subsection*{$L_0$ --- the cited source}

$L_0$ is the ground truth: ranges of Palace's C++ (and the MFEM library beneath it),
quoted by file and line, together with a small set of cross-cutting \emph{reference
notes} that explain the conventions those ranges use---output-argument versus
receiver mutation, the \code{Vector}/\code{ComplexVector} element-type duality,
free-function versus method-form symbols, transparent versus load-bearing
optimization tricks. $L_0$ is not ``ours''; it is the machine as written. Nothing is
interpreted here, only read and cited. It has turned \emph{no} impedance---there is no
layer below it---and every claim anywhere higher in the stack ultimately rests on an
$L_0$ citation. The reference notes are deliberately lean (a few paragraphs of
interpretation plus representative citations each), so that the higher layers can
\emph{point} at a convention rather than re-state it. $L_0$ is where the grep-and-IDE
workflow lives; the citation format is the plain-text \code{path/file.ext:start-end}
the whole book uses.

\subsection*{$L_1$ --- the mutation rotation}

$L_1$ re-expresses each $L_0$ operation as a \emph{pure function over immutable
tensors}. The destination array that $L_0$ overwrote in place becomes a returned
value; \lstinline[style=l4style]|y.Add(alpha, x)| (a mutating member method) becomes
\lstinline[style=l4style]|y_new = axpy(alpha, x, y_old)|, with the $L_0$ destination
buffer gone from the signature. An operator application \lstinline[style=l4style]|A.Mult(x, y)|
that wrote into \code{y} becomes \lstinline[style=l4style]|y = apply_linop(A, x)|, no
destination mentioned. This is the \emph{mutation rotation}, applied operator by
operator: it has turned the impedance \emph{mutation into purity}. $L_1$ is the layer
the book mostly worked in for its operator definitions (Part~IV)---it is the closest
pure-functional layer to the source, its structure still following the source loop,
with workspace scratch buffers erased and only semantically-meaningful aliasing kept.
Reduction-tree non-associativity (the order in which \code{dot} and \code{nrm2} sum)
is preserved here as a \emph{load-bearing} algebraic claim, not quietly discarded.

\subsection*{$L_2$ --- the fusion rotation}

$L_2$ unfolds each $L_1$ function back into a \emph{composition of base algebraic
primitives}, with the performance tricks erased. A fused kernel that computed
$a x + b y + c z$ in one cache-friendly sweep is decomposed into the scaled sums it
stands for; a packed, blocked, SIMD-tuned matrix apply becomes the plain linear map it
implements; a batched specialized BLAS call becomes a composition of base operations.
This is the \emph{fusion rotation}: it has turned the impedance \emph{a fused kernel
into its base algebra}. The discipline of \cref{ch:bigidea} governs the erasure---a
\emph{transparent} performance trick (fusion, tiling, packing) is unfolded and gets a
one-line note, while a \emph{load-bearing} numerical trick (a non-associative reduction
ordering, a deterministic-versus-atomic accumulation) is \emph{preserved} as an
explicit algebraic claim with the property it buys called out. $L_2$ is where the
fold cohorts live: \code{linear\_combination} folds the term axis to a tensor,
\code{inner\_product} folds the length axis to a scalar, and the once-standalone leaves
(\code{axpy}, \code{dot}, \code{scal}, \dots) become thin specializations of their
combinator.

\subsection*{$L_3$ --- the iteration rotation}

$L_3$ lifts the explicit per-degree-of-freedom loops of $L_2$, where the algebra
permits, to \emph{whole-field operations}: ``add these two fields,'' ``apply this
operator to this field,'' with no surviving element loop in the spec form. This is the
\emph{iteration rotation}: it has turned the impedance \emph{an explicit element loop
into a whole-field operation}. It is the rotation that does \emph{not} always go
through---and where it fails, the failure is recorded as a first-class
\emph{obstruction} (\cref{sec:appobstruction}). A clean BLAS-1 update lifts trivially
(every entry independent); a Gauss--Seidel sweep does not (each entry reads the
previous one). $L_3$ therefore arranges its operators on an obstruction-profile
spectrum, from the obstruction-free leaves (\code{apply\_linop}, \code{dot},
\code{normalize}) through operators that carry an obstruction \emph{by reference}
(the divergence-free projector, whose inner solve iterates) to the four
\emph{partial-obstruction} operators whose body lifts cleanly but whose loop does not
(\code{chebyshev}, \code{eigsolve}, \code{orthogonalize}, \code{fold\_solve}).

\subsection*{$L_4$ --- the calculus}

$L_4$ is the small, formally defined graph-evaluation calculus the book taught in
Part~III: high-order combinators (\cref{ch:fold,ch:combinators}), the state monad
(\cref{ch:monad}), the three-stratum state stratification (\cref{ch:strata}), and
operators as values (\cref{ch:operators}). It is \emph{vocabulary, not architecture}---it
says what the algorithm \emph{is}, abstractly, in a form an accelerator backend can
consume, without committing to eager-versus-traced execution or to any particular
container or monad implementation. The \emph{calculus rotation} from $L_3$ replaces
hand-written control flow with named combinators: it has turned the impedance
\emph{hand-written coordination into a named combinator}. A solver's restart-and-converge
loop, threaded by hand at $L_3$, becomes \code{ksp\_solve} folding \code{krylov\_step}
through \code{iterate\_while}; a time march becomes \code{fold\_solve} over a schedule;
an RHS family becomes \code{solve\_family} mapping a fixed operator. $L_4$ is the
backend-lowering target---the whole-tensor, no-aliasing surface the GPU wants---and the
implementation rendered \emph{as} that surface is the synthesized library of
\cref{ch:synthlib}.

\begin{notationbox}
The layer numbers run \emph{opposite} to the direction of construction. The
\emph{stack} is numbered ground-up ($L_0$ source at the bottom, $L_4$ calculus at the
top), but the \emph{rotations} are performed and read \emph{up} the stack: a rotation
``$L_n\!\to\!L_{n+1}$'' takes the lower form to the higher one. A
\emph{lowering}---the spec's name for the formal correspondence---points the other
way, ``$L_{n+1}>L_n$,'' and narrates how the higher form \emph{lowers} into the lower
one. Both describe the same edge; the rotation is the act, the lowering is the
recorded correspondence. The spec stores lowerings in per-adjacent-edge directories
(\code{L4-L3/}, \code{L3-L2/}, \code{L2-L1/}, \code{L1-L0/}), one chapter per theme.
\end{notationbox}

\section{The four rotations between adjacent layers}
\label{sec:approtations}

The four edges of the stack are the four rotations. Each is a \emph{translation across
vocabularies}, not a one-to-one renaming (\cref{def:rotation}): the lower form is
re-expressed in a genuinely different idiom in which one impedance has been turned. A
move that merely pastes new names on the old operations---same state threaded, same
grain, same admitted algorithm in every slot---turns no impedance and is a
\emph{renaming}, the smell \cref{sec:tests} teaches you to catch. \Cref{tab:rotations}
is the reference table: each layer pair, the impedance it rotates, what changes in the
vocabulary, and a worked example drawn from the book.

\begin{table}[t]
\centering
\caption[The four rotations: a reference table]{The four rotations between adjacent
layers, as a reference table. Each row names the edge, the impedance it turns, the
vocabulary shift it performs (the genuine translation, not a rename), and a concrete
example met in the book. Read up the stack ($L_0$ at the bottom): mutation becomes
purity, fusion becomes base algebra, the element loop becomes a whole-field operation,
and hand-written coordination becomes a named combinator.}
\label{tab:rotations}
\small
\renewcommand{\arraystretch}{1.35}
\begin{tabular}{@{}>{\raggedright\arraybackslash}p{0.115\textwidth}
                   >{\raggedright\arraybackslash}p{0.155\textwidth}
                   >{\raggedright\arraybackslash}p{0.30\textwidth}
                   >{\raggedright\arraybackslash}p{0.30\textwidth}@{}}
\toprule
\textbf{Edge} & \textbf{Impedance turned} & \textbf{What rotates (the translation)} &
\textbf{Example} \\
\midrule
\textbf{$L_4 > L_3$} \newline calculus rotation &
hand-written control flow $\to$ named combinator &
The monad / strata / \code{readonly} wrappers dissolve; coordination loops become
\code{iterate\_while}, \code{fold\_solve}, \code{solve\_family}. State that was
threaded by hand is hidden in a combinator's contract. &
\code{ksp\_solve} folding \code{krylov\_step} through \code{iterate\_while} replaces a
hand-threaded restart loop; the time march becomes \code{fold\_solve} over a schedule
(\cref{ch:transient}). \\
\textbf{$L_3 > L_2$} \newline iteration rotation &
explicit element loop $\to$ whole-field op &
The per-DoF loop is erased and the operation re-stated on whole tensors---\emph{or},
where it cannot be, an \code{obstruction} marker is named first-class and shadows to an
$L_2$ non-law. &
\lstinline[style=l4style]|for i: y[i]=a*x[i]+y[i]| becomes the single whole-field
\code{axpy}; the Gauss--Seidel sweep refuses and records an obstruction
(\cref{sec:appobstruction}). \\
\textbf{$L_2 > L_1$} \newline fusion rotation &
fused kernel $\to$ base algebra &
A fused / packed / batched kernel is unfolded into a composition of base primitives;
transparent tricks get a note, load-bearing numerical tricks are preserved as explicit
claims. &
A fused $ax+by+cz$ sweep unfolds into scaled sums; the leaves \code{axpy}/\code{dot}
become specializations of the \code{linear\_combination} / \code{inner\_product}
folds. \\
\textbf{$L_1 > L_0$} \newline mutation rotation &
mutation $\to$ purity &
The in-place destination buffer disappears from the signature; a mutating method
becomes a function returning a fresh value. The lowering reintroduces the buffer when
reading down to source. &
\lstinline[style=l4style]|y.Add(alpha, x)| (writes \code{y}) becomes
\lstinline[style=l4style]|axpy alpha x y| (returns fresh);
\lstinline[style=l4style]|A.Mult(x, y)| becomes \code{apply\_linop}
(\cref{ch:bigidea}). \\
\bottomrule
\end{tabular}
\end{table}

A few structural facts about the table are worth stating once. First, not every edge
is substantive for every operator: a leaf that is already at the right grain---an
\code{axpy} that is whole-field by construction---``carries through'' an edge unchanged,
and the spec records that as an in-line identity note rather than a separate lowering
theme (a mirrored entry with a thin theme would be the renaming smell). The genuine
content of an edge lives in the operators that \emph{do} rotate: the solver caps at
$L_4>L_3$, the fold cohorts at $L_2>L_1$, the constructed-operator gates and the
sequential obstructions at $L_3>L_2$.

Second, an edge can be substantive in one direction's vocabulary and identity-in-form
in the other. The Krylov step's \emph{body} is identity-in-form across $L_4>L_3>L_2$
(the same primitive sequence, only the wrappers dissolving), while its \emph{driver}---the
outer restart-and-converge loop---rotates substantively at every edge (the iteration
view is rendered at $L_3$, erased at $L_2$, named as a combinator at $L_4$). The spec
captures this as a kernel-identity-plus-driver-non-identity pair of themes.

Third, identity rotations that span \emph{non-adjacent} layers are annotated in-line,
not given their own directory. There is no \code{L3-L1/}; when an operator's $L_3$ body
is value-thread-isomorphic to its $L_1$ form through an identity-like $L_2$ absorption,
that relationship is the transitive composition of the adjacent-edge themes, noted in
the $L_n$ entry.

\section{Obstructions: the things that do not rotate}
\label{sec:appobstruction}

Not everything rotates. Some work, faithfully described at $L_2$, genuinely cannot be
lifted to a clean $L_3$ whole-field operation, because its very meaning depends on doing
things in order. When that happens the spec does \emph{not} fake a rotation and does
\emph{not} treat the failure as a defect in the method. It records an
\emph{obstruction}: a precise statement of what blocks the lift and why. An obstruction
is a first-class negative result---a finding, not a failure (\cref{ch:rotations}). The
obstructions are some of the most valuable output in the stack, because they say
exactly where the abstract whole-tensor picture meets a wall.

There are two \emph{kinds} of obstruction, distinguished by how the blocking loop
relates to Palace, and the kind determines whether the obstruction can ever be
discharged. \Cref{tab:obstructions} catalogues both, with the book chapters where each
was met; \cref{fig:obstructionkinds} draws the distinction.

\begin{table}[t]
\centering
\caption[The obstruction catalogue]{The catalogue of non-rotation. Two kinds of
obstruction block the iteration rotation, distinguished by \emph{whose} loop carries
the blockage. A \emph{sequential-obstruction} is a Palace-authored loop whose
$i$-th step reads the $(i\!-\!1)$-th step's output (a loop-carried dependency); Palace
\emph{renders} it, so it survives as ``these reordering laws do not hold,'' and it
could in principle be re-described. An \emph{opaque-library-ownership} obstruction is a
loop that lives entirely in third-party code (SLEPc, MFEM); Palace authors no loop to
render, only a boundary marker, and it cannot be discharged unless the upstream
architecture changes.}
\label{tab:obstructions}
\small
\renewcommand{\arraystretch}{1.3}
\begin{tabular}{@{}>{\raggedright\arraybackslash}p{0.20\textwidth}
                   >{\raggedright\arraybackslash}p{0.34\textwidth}
                   >{\raggedright\arraybackslash}p{0.34\textwidth}@{}}
\toprule
\textbf{Kind} & \textbf{What blocks the lift} & \textbf{Where the book met it} \\
\midrule
\textbf{sequential-} \newline \textbf{obstruction} \newline \emph{(Palace-authored
loop)} &
A loop-carried dependency: step $i$ reads step $i\!-\!1$'s output, so there is no
``all entries at once'' form. Palace \emph{renders} the loop (an explicit tail
recursion); the lift erases the iteration view and the obstruction survives as a lost
reordering law. Root is usually numerical stability. &
The MGS arm of orthogonalization---the serial rank-1-projector chain that buys
roundoff-orthogonality (\cref{ch:orthog}); the Chebyshev smoother's inner $k$-recurrence
and outer Richardson sweep; the Gauss--Seidel-flavored multigrid smoother sweep; the
transient time march's carry-threading (\cref{ch:transient}). \\
\textbf{opaque-} \newline \textbf{library-} \newline \textbf{ownership} \newline
\emph{(third-party loop)} &
The iteration lives entirely inside a library Palace calls but does not implement.
Palace authors \emph{no} loop, so there is nothing to render---only a marker at the
library boundary. No promotion route unless Palace re-architects; these are the
\code{\#extern} edges of \cref{app:kernel}. &
The SLEPc / ARPACK eigensolve loop, handed the operators and run to convergence inside
the library (\cref{ch:eigenmodes}); the MFEM ODE time-integrator step inside
\code{fold\_solve} (\cref{ch:transient}); the libCEED element-quadrature kernel inside
\code{fe\_assemble}. \\
\bottomrule
\end{tabular}
\end{table}

\begin{figure}[t]
\centering
\begin{tikzpicture}[node distance=6mm]
  % ---- LEFT: liftable, independent ----
  \begin{scope}
    \node[font=\bfseries\small, simGreen] at (1.4,1.85) {Liftable: independent updates};
    \foreach \i in {0,1,2,3}{
      \node[opbox, minimum width=8mm, minimum height=6mm] (a\i) at (\i*0.95,0.7) {$y_\i$};
      \node[data, minimum width=8mm, minimum height=5mm] (x\i) at (\i*0.95,-0.4) {$x_\i$};
      \draw[flow] (x\i) -- (a\i);
    }
    \node[align=center, font=\scriptsize\itshape, simGray, text width=42mm] at (1.4,-1.3)
      {each $y_i = a x_i + y_i$ reads only its own input $\Rightarrow$ one whole-field op};
    \node[draw=simGreen, thick, rounded corners=2pt, fit=(a0)(a3), inner sep=3pt,
          fill=simGreen!8] {};
    \node[font=\scriptsize\bfseries, simGreen] at (1.4,1.2) {$L_3$ whole-field op};
  \end{scope}
  % ---- MIDDLE: sequential-obstruction (Palace-authored) ----
  \begin{scope}[shift={(5.4,0)}]
    \node[font=\bfseries\small, simRust, align=center, text width=42mm] at (1.4,1.85)
      {Sequential-obstruction\\[-1pt]\scriptsize(Palace authors the loop)};
    \foreach \i in {0,1,2,3}{
      \node[opbox, draw=simRust, minimum width=8mm, minimum height=6mm] (b\i) at (\i*0.95,0.7) {$y_\i$};
      \node[data, minimum width=8mm, minimum height=5mm] (u\i) at (\i*0.95,-0.4) {$x_\i$};
      \draw[flow] (u\i) -- (b\i);
    }
    \foreach \i/\j in {0/1,1/2,2/3}{
      \draw[backflow] (b\i.east) -- (b\j.west);}
    \node[align=center, font=\scriptsize\itshape, simGray, text width=42mm] at (1.4,-1.3)
      {$y_i$ reads the \emph{already-updated} $y_{i-1}$ $\Rightarrow$ loop renders, does not lift};
    \node[font=\scriptsize\bfseries, simRust] at (1.4,1.2) {loop-carried dependency};
  \end{scope}
  % ---- RIGHT: opaque-library-ownership ----
  \begin{scope}[shift={(10.8,0)}]
    \node[font=\bfseries\small, simViolet, align=center, text width=40mm] at (1.1,1.85)
      {Opaque-library-ownership\\[-1pt]\scriptsize(Palace authors no loop)};
    \node[extern, minimum width=34mm, minimum height=18mm, align=center] (lib) at (1.1,0.15)
      {\textbf{\#extern}\\[2pt]\scriptsize the loop lives entirely\\in SLEPc / MFEM / libCEED;\\only a boundary marker};
    \node[data, minimum width=10mm, above=4mm of lib] (in2) {operators in};
    \node[data, minimum width=10mm, below=4mm of lib] (out2) {result out};
    \draw[flow] (in2) -- (lib);
    \draw[flow] (lib) -- (out2);
    \node[align=center, font=\scriptsize\itshape, simGray, text width=40mm] at (1.1,-1.95)
      {nothing to render $\Rightarrow$ a marker, not a recurrence};
  \end{scope}
\end{tikzpicture}
\caption[The two kinds of obstruction]{Why some loops rotate to $L_3$ and some do not,
and the two kinds of those that do not. \emph{Left (green):} an \code{axpy}-style
update writes $y_i = a x_i + y_i$ reading only its own input; the entries are
independent, so the loop lifts to one whole-field operation. \emph{Middle (rust):} a
\emph{sequential-obstruction}---a Gauss--Seidel-like sweep where $y_i$ reads the
already-updated $y_{i-1}$ (the dashed loop-carried dependency). Palace authors this
loop, so $L_3$ \emph{renders} it as an explicit tail recursion; the lift to $L_2$
erases the iteration view and the obstruction survives as a lost reordering law. It
could, in principle, be re-described. \emph{Right (violet):} an
\emph{opaque-library-ownership} obstruction---the loop lives entirely in third-party
code (the eigensolver, the ODE step, the quadrature kernel). Palace authors no loop, so
there is nothing to render; only a boundary marker, the \code{\#extern} edge of
\cref{app:kernel}, and no discharge route unless the upstream changes.}
\label{fig:obstructionkinds}
\end{figure}

The distinction is not pedantry; it changes how an obstruction can ever be removed. A
\emph{rendered} recurrence (the sequential kind) records a loop Palace chose to
surface and could in principle re-describe---and indeed, for the three opaque kernels
that are spine dependencies, the book provides \emph{both} an opaque
\code{\#extern} contract and a separate, constructive realization in our own
vocabulary, linked by a reviewable correspondence (\cref{app:kernel}'s
kernel-API/kernel-implementation split; the eigensolve's Krylov--Schur reconstruction,
\cref{ch:eigenmodes}). An \emph{opaque-library} marker, by contrast, records a boundary
Palace never sees inside and will not discharge unless the upstream architecture
changes.

A finer reading---the \emph{erasure scope} of \cref{ch:rotations}---asks how much of
the iteration view a single layer-hop erases: a single sequential loop (a linear
solve's outer iteration), a nested double loop (a polynomial smoother's inner
recurrence inside an outer sweep), or an opaque library call with no loop to render at
all. The rung names what kind of obstruction you are looking at; \cref{ch:rotations}'s
erasure-scope ladder draws all three.

\begin{pitfall}
An obstruction is never an excuse to stop, and never a whole-operator verdict when
only the loop resists. The honest record is precise about the boundary: the per-step
\emph{body} usually rotates cleanly (it is identity-in-form to the layer below), and
only the \emph{iteration} fails. That is exactly the \emph{partial-obstruction} status
the spec uses for \code{chebyshev}, \code{eigsolve}, \code{orthogonalize}, and
\code{fold\_solve}---``the body lifts; the loop does not, because of this cited
dependency.'' Reading one of these as a total failure, or forcing a false whole-field
form to avoid recording the obstruction, both throw away the finding the obstruction
\emph{is}.
\end{pitfall}

\section{How the book navigated the stack}
\label{sec:appnavigation}

The book did not present the five layers in order. It taught the \emph{top}---the $L_4$
calculus---as the main object (Part~III), grounded it in the \emph{bottom}---the cited
source and the pure-function operators (Parts~II and~IV)---and then \emph{composed} the
top into the user-facing modalities (Part~V). The synthesized library (Part~VI)
rendered that same $L_4$ surface as code. \Cref{fig:partsmap} charts which part teaches
which layer, so the reader who now wants the full spec knows where each piece of the
book already sits in the stack.

\begin{figure}[t]
\centering
\begin{tikzpicture}[font=\footnotesize, >=Stealth, node distance=0mm]
  \def\lw{58mm}
  \def\lh{9mm}
  \def\gap{3.0mm}
  % the five layer bands (left column)
  \foreach \i/\name/\col/\bg in {%
      0/$L_0$ cited source/simInk/simBgGray,
      1/$L_1$ pure functions/simGreen/simBgGreen,
      2/$L_2$ algebra/simTeal/simBgGray,
      3/$L_3$ tensor-field ops/simBlue/simBgBlue,
      4/$L_4$ the calculus/simViolet/simBgViolet}{
    \node[draw=\col, fill=\bg, thick, rounded corners=2pt,
          minimum width=\lw, minimum height=\lh, align=center, font=\small]
      (B\i) at (0, \i*\dimexpr\lh+\gap\relax) {\textbf{\name}};
  }
  % the part annotations (right column), aligned to layers, with connector braces
  \def\px{0.5*\lw+6mm}
  % Parts II/IV -> L0/L1  (grounding)
  \draw[decorate, decoration={brace, amplitude=5pt, mirror}, simGreen, thick]
    ($(B0.east)+(3mm,-0.5*\lh+1mm)$) -- ($(B1.east)+(3mm,0.5*\lh-1mm)$);
  \node[product, anchor=west, text width=44mm, align=left] at (\px, 0.5*\dimexpr\lh+\gap\relax)
    {\textbf{Parts II \& IV} $\to$ $L_0$/$L_1$\\[1pt]\scriptsize the field theory + the
     pure-function operators: the bottom-up grounding};
  % L2 -> intervening
  \node[opbox, anchor=west, text width=44mm, align=left] at (\px, 2*\dimexpr\lh+\gap\relax)
    {\textbf{the intervening algebra}\\[1pt]\scriptsize $L_2$/$L_3$: the stepping-stones
     the rotations pass through (the spec's middle)};
  \draw[refedge] (B2.east) -- (\px-1mm, 2*\dimexpr\lh+\gap\relax);
  % Part III -> L4 (the taught surface)
  \draw[decorate, decoration={brace, amplitude=5pt, mirror}, simViolet, thick]
    ($(B4.east)+(3mm,-0.5*\lh+1mm)$) -- ($(B4.east)+(3mm,0.5*\lh-1mm)$);
  \node[stage, fill=simBgViolet, draw=simViolet, very thick, anchor=west,
        text width=44mm, align=left] at (\px, 4*\dimexpr\lh+\gap\relax)
    {\textbf{Part III} $\to$ $L_4$ \emph{(the taught surface)}\\[1pt]\scriptsize the calculus:
     fold, combinators, monad, strata, operators-as-values};
  % Part V / VI annotations along the top
  \node[font=\scriptsize\itshape, simRust, align=center, text width=0.95\linewidth]
    at (0.30\linewidth, 4*\dimexpr\lh+\gap\relax+0.5*\lh+5mm)
    {\textbf{Part V} reads the modalities as $L_4$ compositions (drivers as composition
     roots); \quad \textbf{Part VI} renders the $L_4$ surface as the synthesized library
     (\cref{ch:synthlib}).};
\end{tikzpicture}
\caption[How the book's parts map onto the layers]{The correspondence between the
book's parts and the stack's layers. The book taught the \emph{top} of the stack as its
primary object---\textbf{Part~III} is the $L_4$ calculus (violet, the taught
surface)---and grounded it in the \emph{bottom}: \textbf{Parts~II and~IV} supply the
$L_0$ cited source and the $L_1$ pure-function operators (green, the bottom-up
grounding). The intervening $L_2$/$L_3$ algebra (teal/blue) is the spec's middle---the
stepping-stones the rotations pass through, which the book touched mainly when an
operator's iteration rotation or fusion rotation was load-bearing. \textbf{Part~V} then
reads the five modalities as $L_4$ compositions (each driver a composition root over
the calculus), and \textbf{Part~VI} renders that same $L_4$ surface as the synthesized
library. The reader who wants the full specification now has the map: the book is the
$L_4$-and-grounding view; the spec is all five layers and the rotations between them.}
\label{fig:partsmap}
\end{figure}

The choice to teach $L_4$ first, rather than climb the stack from $L_0$, was
deliberate. A direct $L_0\!\to\!L_4$ leap would be a leap of faith; the stack replaces
it with a sequence of short, auditable rotations, but a \emph{reader} does not need to
re-perform that sequence to understand the result. The result---what the algorithm
\emph{is}, abstractly---is the $L_4$ calculus, and that is what the book taught. The
grounding in Parts~II and~IV is what lets you trust the calculus is faithful; the
rotations of \cref{ch:rotations} are what connect the two; and this appendix is the map
that lets you go read the rest. The payoff of the whole layered construction is three
things at once: \emph{auditability} (each short move is checkable), \emph{honesty} (the
obstructions are recorded, not papered over), and a \emph{target} (the $L_4$ surface an
accelerator backend consumes, rendered as the library of \cref{ch:synthlib}).

\summarysection
The simulator is described not once but at a stack of five \emph{layers}, each complete
and correct in its own vocabulary: $L_0$, the cited Palace/MFEM source plus its
reference notes; $L_1$, pure functions over immutable tensors; $L_2$, algebraic
decompositions with HPC fusions erased; $L_3$, global tensor-field operations; and
$L_4$, the calculus of Part~III. Adjacent layers are joined by one of \emph{four
rotations}, each turning exactly one impedance while computing the same numbers: the
\emph{mutation rotation} ($L_1>L_0$, mutation into purity), the \emph{fusion rotation}
($L_2>L_1$, a fused kernel into base algebra), the \emph{iteration rotation} ($L_3>L_2$,
an element loop into a whole-field operation), and the \emph{calculus rotation}
($L_4>L_3$, hand-written control flow into a named combinator). Every rotation is a
translation across vocabularies, never a renaming. Where the iteration rotation cannot
go through, the spec records a first-class \emph{obstruction}---of two kinds: a
\emph{sequential-obstruction} (a Palace-authored loop-carried dependency, rendered and
re-describable in principle---Gauss--Seidel smoothing, MGS orthogonalization, the
Chebyshev recurrence) and an \emph{opaque-library-ownership} obstruction (a loop owned
by SLEPc, MFEM, or libCEED, marked but not rendered, the \code{\#extern} edges of
\cref{app:kernel})---each tied to the chapter where the book met it
(\cref{ch:orthog,ch:eigenmodes,ch:transient}). The book navigated the stack by teaching
the top ($L_4$, Part~III), grounding it in the bottom ($L_0$/$L_1$, Parts~II and~IV),
and composing it into the modalities (Part~V) and the synthesized library (Part~VI).
The reader who wants the full specification now holds the map: five layers, four
rotations, and the obstructions that, by not rotating, are findings rather than
failures.

\furtherreading
The layered, rotation-based method this appendix maps is purpose-built for this project,
but the obstructions it catalogues are the standard hard cases of numerical linear
algebra. The sequential reductions that resist the iteration rotation---Gauss--Seidel
and SSOR smoothing, the modified-Gram--Schmidt orthogonalization whose serial
rank-one projector chain buys the roundoff-orthogonality the batched form loses---are
developed in Saad's \emph{Iterative Methods for Sparse Linear
Systems}~\citep{saad2003}, the standard reference for the Krylov and multigrid loops
behind the single-loop and nested-loop obstructions of \cref{sec:appobstruction}.
Trefethen and Bau's \emph{Numerical Linear Algebra}~\citep{trefethenbau1997} develops
orthogonalization, conditioning, and the eigenvalue problem lecture by lecture, and is
the ideal companion for seeing \emph{why} the eigensolve loop is the canonical
opaque-library obstruction---why the iteration is best left to a specialized library
and marked at the boundary rather than re-derived. The simulator whose five-layer
dissection this map summarizes is Palace~\citep{palace2023}; reading its source against
this stack, with the four rotations of \cref{sec:approtations} and the obstruction
catalogue of \cref{sec:appobstruction} in hand, is exactly the auditable
ground-to-calculus walk the layered method was built to make possible.
